\documentclass[10pt,a4paper]{amsart}
\usepackage[margin=19mm]{geometry}
\usepackage{amsmath,amssymb,mathtools}
\usepackage[scr=rsfs,cal=euler]{mathalfa}
\usepackage{xfrac}
\usepackage[unicode,hidelinks,bookmarks=false]{hyperref}
\newcommand{\ie}{i.e.\ }
\newcommand{\cf}{cf.\ }
\newcommand{\resp}{respectively\ }
\newcommand{\Subsection}{\subsection}
\providecommand{\nobreakdash}{\nobreak\hbox{-}}
\setlength{\emergencystretch}{2em}
\multlinegap0pt
\usepackage{bm}
\usepackage{bbm}
\usepackage{dsfont}
\usepackage{xspace}
\usepackage{cancel}
\usepackage{mathtools}
\usepackage{amsthm}
\makeatletter
\@ifundefined{definition}{\newtheorem{definition}{Definition}}{}
\@ifundefined{lemma}{\newtheorem{lemma}{Lemma}}{}
\@ifundefined{proposition}{\newtheorem{proposition}{Proposition}}{}
\@ifundefined{theorem}{\newtheorem{theorem}{Theorem}}{}
\makeatother
\makeatletter
\expandafter\ifx\csname example\endcsname\relax
\newenvironment{example}[1][]%
    {\begin{innerexample}[#1]\pushQED{\qed}}%
    {\popQED\end{innerexample}}
\fi
\expandafter\ifx\csname proof\endcsname\relax
\renewenvironment{proof}[1][Proof]{\par\pushQED{\qed}\normalfont\textit{#1.}~}{\popQED\par}
\fi
\makeatother
\title[Modular Design of Strict Control Lyapunov Functions for Glob]{Modular Design of Strict Control Lyapunov Functions for Global Stabilization of the Unicycle in Polar Coordinates}
\author{Velimir Todorovski, Kwang Hak Kim,, Miroslav Krstic}
\date{Source: arXiv:2509.25575v1 (CC BY 4.0)}
\begin{document}
\maketitle

\noindent\emph{Source and licence.} Modular Design of Strict Control Lyapunov Functions for Global Stabilization of the Unicycle in Polar Coordinates. Version arXiv:2509.25575v1 (CC BY 4.0), distributed under the Creative Commons Attribution 4.0 International licence, as recorded in the arXiv licence metadata for this exact version and shown on the abstract page https://arxiv.org/abs/2509.25575v1. Full rights notice: licenses/arxiv-2509.25575-CC-BY-4.0.txt.\par\smallskip

\noindent\textit{Editorial scope.} Counted unit: the Theorem whose source label is \texttt{thm:unicycle\_CLF\_BoLSA} in the article (source0.tex lines 567--593, source label \texttt{thm:unicycle\_CLF\_BoLSA}), delivered together with the article's own complete original proof of it (source0.tex lines 595--693, one \texttt{proof} environment of 99 lines). No mathematics is written, repaired or re-derived: statement and proof are the article's own text line for line. Required context retained uncounted: eq:unicycle\_polar\_closed\_loop-Gv-1 (lines 174--182); eq-basic-v-control (lines 218--221); eq:unicycle\_polar\_closed\_loop-Gv-2 (lines 224--228); eq-omega-general (lines 240--244); eq:unicycle\_polar\_closed\_loop-Gv-3-pass (lines 246--253); def-our-GAS (lines 345--352); prop:composite\_Lyap\_function (lines 362--413); def-CLF (lines 415--430); eq:backstepping\_z (lines 471--474); eq:psi\_definition (lines 476--483); thm:backstepping\_theorem (lines 487--507); eq:Delta\_Barfli (lines 495--498); eq:Delta\_globa (lines 495--498); lemma:sinterm\_upperbound (lines 561--563). Every definition, display and label of those lines is retained unchanged. Omitted from this package and not redistributed: 1--173, 183--217, 222--223, 229--239, 245--245, 254--344, 353--361, 414--414, 431--470, 475--475, 484--486, 499--560, 564--566, 594--594, 694--1004. Nothing inside a retained excerpt is omitted or altered: every retained body is delivered exactly as the article's lines are written, so no omission receipt of any kind accompanies this package. Citations: the retained excerpts carry no citation command at all, so no bibliography entry of the article is transcribed and no reference is redistributed. Reference adaptation: none was needed and none was made; every reference inside the delivered unit resolves inside it, so no unresolved reference or citation is produced. Printed slips kept literal: no printed word, symbol, hypothesis or number of the counted statement or of its proof was altered. Layout and class: the article's own class is \texttt{ieeeconf} with packages including mathrsfs, xcolor, amsmath, amssymb, amsfonts, graphicx, algorithm, algorithmic, textcomp, setspace, lipsum, mathtools, enumerate, caption; the wrapper is the lane's pinned \texttt{amsart} editorial wrapper at 10pt on A4, which restates the theorem-like environments the retained text needs and adds the article's own macro definitions that the retained text needs (none), with extra wrapper packages (bm, bbm, dsfont, xspace, cancel, mathtools). The delivered numbering is the unit's own. Assets: no retained span contains a figure environment, a graphics-inclusion command, a PDF-page inclusion, a table or a TikZ picture; no image file is copied into this package, the package declares no asset dependency and no image file is redistributed with it. Licence: version arXiv:2509.25575v1 of this article is distributed under the Creative Commons Attribution 4.0 International licence, as recorded in the arXiv licence metadata for this exact version and shown on the abstract page https://arxiv.org/abs/2509.25575v1. A full rights notice is delivered at licenses/arxiv-2509.25575-CC-BY-4.0.txt. Characters: every retained excerpt is pure ASCII, so the delivered engine, XeLaTeX with its default Latin Modern text font, reports no missing character.
\begin{subequations}\label{eq:unicycle_polar_closed_loop-Gv-1}
    \begin{align}
    \dot{\rho} &= 
    %- k_1 \rho \cos^2(\gamma)
    -v \cos\gamma\label{eq:unicycle_polar_rhodot}\\
    \dot{\delta} &= v  \frac{\sin\gamma}{\rho}\label{eq:unicycle_polar_deltadot}\\
    \dot{\gamma} &= v \frac{\sin\gamma}{\rho} -\omega \label{eq:unicycle_polar_gammadot} \,.
    \end{align} 
\end{subequations}

\begin{equation}
\label{eq-basic-v-control}
\fbox{$v =  k_1 \rho \cos\gamma$} = -k_1(x\cos\theta+y\sin\theta)\,,  \quad k_1>0\,.
\end{equation}

\begin{align}
\label{eq:unicycle_polar_closed_loop-Gv-2}
\dot{\rho} &=  
- k_1 \rho \cos^2(\gamma)\,,
\end{align} 

\begin{equation}
\label{eq-omega-general}
%, \eqref{eq:angular_velocity_genova}
\fbox{$\omega = \dfrac{k_1}{2} \sin(2\gamma) +\tilde\omega$}
\end{equation}

\begin{subequations}
\label{eq:unicycle_polar_closed_loop-Gv-3-pass}
\begin{eqnarray}
\label{eq:unicycle_polar_closed_loop-Gv-3n-pass}
\dot{\delta} &=&  \frac{k_1}{2} \sin(2\gamma)\\ \label{eq:unicycle_polar_closed_loop-Gv-3a-pass}
\dot{\gamma} &=& -\tilde\omega\,.
\end{eqnarray}
\end{subequations}

\begin{definition}
\label{def-our-GAS}
Consider the system \eqref{eq:unicycle_polar_closed_loop-Gv-2}, \eqref{eq:unicycle_polar_closed_loop-Gv-3-pass} with a feedback law $\tilde\omega(\gamma,\delta)$ that is continuous on a state space $\mathcal{Q}$ with respect to its metric. 
%on the state space $\mathcal{S}_i$ with $i = 0,1,2,3$. 
If there exists a class $\mathcal{KL}$ function $\beta$ such that, for all $t\geq 0$, it holds that $|(\rho(t),\delta(t), \gamma(t))|_{\mathcal{Q}}\leq \beta\left(|(\rho_0,\delta_0, \gamma_0)|_{\mathcal{Q}},t\right)$, we say that the 
%equirlibrium 
point $\rho=\delta = \gamma=0$ is {\em globally asymptotically stable on $\mathcal{Q}$} (GAS on $\mathcal{Q}$). 
\end{definition}

\begin{proposition}[{Composite Lyapunov functions}] 
\label{prop:composite_Lyap_function}
Consider 
%any of the example functions $V_\rho$ in \eqref{eq-rhoCLFs} and 
any continuously differentiable function $(\delta,\gamma) \mapsto V_{\delta\gamma}$ where $(\delta,\gamma)$ belong to any state space $\hat{\mathcal{T}}\in\{\mathcal{T},\mathcal{T}_1,\mathcal{T}_2,\mathcal{T}_3\}$
and  such that 
$ \alpha_1(|(\delta, \gamma)|_{\hat{\mathcal{T}}}) \le V_{\delta\gamma}(\delta,\gamma) \le \alpha_2(|(\delta, \gamma)|_{\hat{\mathcal{T}}})$ where $\alpha_1, \alpha_2$ are class $\mathcal{K}_{\infty}$ functions. 
Let $(\delta,\gamma) \mapsto\tilde\omega$ be a continuous function such that 
% \begin{align}
% \dot V_{\delta\gamma}(\delta,\gamma) &:= \dfrac{\partial V_{\delta\gamma}}{\partial\delta}
% \frac{k_1\sin(2\gamma)}{2}  -\dfrac{\partial V_{\delta\gamma}}{\partial\gamma}\tilde\omega(\delta,\gamma) \nonumber \\
% &\leq -\alpha_{\delta\gamma}(V_{\delta\gamma}) \label{eq:V_dot_delta_gamma}
% \end{align}
\begin{equation}
\dot V_{\delta\gamma} := \dfrac{\partial V_{\delta\gamma}}{\partial\delta}
\frac{k_1\sin(2\gamma)}{2}  -\dfrac{\partial V_{\delta\gamma}}{\partial\gamma}\tilde\omega(\delta,\gamma) \leq -\alpha_{\delta\gamma}(V_{\delta\gamma}) \label{eq:V_dot_delta_gamma}
\end{equation}
for some class $\mathcal{K}$ function $\alpha_{\delta\gamma}$. 
%whenever $\rho>0$ or $\delta\neq 0$ or $\gamma\neq 0$. 
Then for {\em any} function $(r,s) \mapsto \mathcal{V}$ such that 
\begin{enumerate}
\item $\mathcal{V}(0,0)=0$ and $\mathcal{V}(r, s)>0$ if $r > 0 $ or $s > 0$, \label{it:pos_def_composite_Lyap_function} 
\item $\lim_{r+s\rightarrow\infty} \mathcal{V}(r,s) = \infty$,
\item $\dfrac{\partial \mathcal{V} }{\partial r }(r,s)>0$ and $\dfrac{\partial \mathcal{V} }{\partial s}(r,s)>0$  if $r>0$ or $s > 0$, \label{it:positive_partials}
\end{enumerate}
including the particular functions
 $\mathcal{V}(r,s)= r+s$, $\mathcal{V}(r,s)= \ln(1+r) + s$, $\mathcal{V}(r,s)= (1+r){\rm e}^{s} -1$,  
% \begin{enumerate}[(i)]
% \item $\mathcal{V}(r,s)= r+s$ \label{it:(i)}
% \item $\mathcal{V}(r,s)= \ln(1+r) + s$
% \item $\mathcal{V}(r,s)= {\rm e}^{r} - 1 + s$
% \item $\mathcal{V}(r,s)= (1+r){\rm e}^{s} -1$
% \item $\mathcal{V}(r,s) = r+s+rs$ 
% \item $\mathcal{V}(r,s) = \cosh(r) + s - 1 $ \label{it:(v)}
% \end{enumerate}
for both of the `composite' Lyapunov functions $V(\rho,\delta,\gamma) = \mathcal{V}\left(\rho^2,V_{\delta\gamma}(\delta,\gamma)\right)$ and $V(\rho,\delta,\gamma) = \mathcal{V}\left(V_{\delta\gamma}(\delta,\gamma), \rho^2 \right)$ there exist the respective (distinct) 
triplets of functions $\bar\alpha_1,\bar\alpha_2\in\mathcal{K}_\infty$, $\alpha\in\mathcal{K}$
%class $\mathcal{K}$ functions triples $(\bar\alpha_1,\bar\alpha_2,\alpha)$ 
such that,  for all $(\rho,\delta,\gamma)$ in the state space $\hat{\mathcal{S}}\in\{\mathcal{S},\mathcal{S}_1,\mathcal{S}_2,\mathcal{S}_3\}$, it holds that  
$ \bar\alpha_1(|(\rho,\delta, \gamma)|_{\hat{\mathcal{S}}}) \le V(\rho,\delta,\gamma) \le \bar\alpha_2(|(\rho,\delta, \gamma)|_{\hat{\mathcal{S}}})$ and 
\begin{align}
\dot{V}(\rho,\delta,\gamma) := 
-\dfrac{\partial V }{\partial \rho}%(\rho,\delta,\gamma)
k_1\rho^2&\cos^2\gamma  + %\dfrac{\partial V}{\partial s} (\rho,\delta,\gamma)\left[
\dfrac{\partial V}{\partial \delta}
\frac{k_1 \sin(2\gamma)}{2} 
 \nonumber\\
&-\dfrac{\partial V}{\partial\gamma}\tilde\omega(\delta,\gamma) \leq -\alpha(V)\,. \label{eq:V_dot_rho_delta_gamma}
\end{align}

%for all $\rho>0$.
\end{proposition}

\begin{definition}[CLF for the unicycle \eqref{eq:unicycle_polar_closed_loop-Gv-1}]
\label{def-CLF}
A continuously differentiable function $(\rho,\delta,\gamma)\mapsto V$ 
is a \textit{control Lyapunov function} (CLF) for \eqref{eq:unicycle_polar_closed_loop-Gv-1} if it has the following properties.
\begin{enumerate}
\item There exist class $\mathcal{K}$ functions  $(\bar\alpha_1,\bar\alpha_2)$ such that,  for all $(\rho,\delta,\gamma)$ in 
$\Sigma = \{\rho\geq 0\}\times \hat{\mathcal{T}}$, where $\hat{\mathcal{T}}\in\{\mathcal{T},\mathcal{T}_1\}$, 
$ \bar\alpha_1(|(\rho,\delta, \gamma)|_{\hat{\mathcal{S}}}) \le V(\rho,\delta,\gamma) \le \bar\alpha_2(|(\rho,\delta, \gamma)|_{\hat{\mathcal{S}}})$, where $\hat{\mathcal{S}} = \{\rho> 0\}\times \hat{\mathcal{T}}$. \label{it:clf_prop_1}
\item There exists $(v/\rho,\omega)\in \mathbb{R}^2$ such that
\begin{equation*}
\left[-\dfrac{\partial V}{\partial\rho}\rho\cos\gamma + \left(\dfrac{\partial V}{\partial\delta}+\dfrac{\partial V}{\partial\gamma}\right)\sin\gamma \right]\dfrac{v}{\rho} - \dfrac{\partial V}{\partial\gamma}\omega < 0 \,,
\end{equation*}
for all $(\rho,\delta,\gamma)\neq (0,0,0)$ in $\Sigma$. \label{it:clf_prop_2}
\end{enumerate}
% , is referred to as a {\em control Lyapunov function} (CLF) for \eqref{eq:unicycle_polar_closed_loop-Gv-1}.
\end{definition}

\begin{align}
\label{eq:backstepping_z}
\fbox{$\displaystyle z = \gamma + \frac{1}{2}\arctan(2 k_2 \Delta)$}
\end{align}

\begin{align}
&\psi(z,\gamma) = \frac{\sin(2z-2\gamma) +\sin(2\gamma)}{2z}
\nonumber\\
=& \frac{1}{\sqrt{1 + 4k_2^2 \Delta^2}} \left(\,\frac{\sin(2z)}{2z}
% \frac{\sin(2z)}{2z}
+ 2k_2 \Delta \frac{1 - \cos(2z)}{2z}\right)
 \label{eq:psi_definition}
\end{align}

\begin{theorem}
\label{thm:backstepping_theorem}
Consider the system \eqref{eq:unicycle_polar_closed_loop-Gv-1} in closed-loop with \eqref{eq-basic-v-control}, \eqref{eq-omega-general}, and 
\begin{equation}
\tilde\omega = k_4 z + \dfrac{\partial \Delta}{\partial \delta} \left( \frac{k_1k_2 \sin(2\gamma)}{2(1+4k_2^2 \Delta^2)} + k_3 \psi(z,\gamma)\Delta \right),
\label{eq:backstepping_controllers}
\end{equation}
where $k_1,k_2,k_3,k_4 > 0$, and $z$ and $\psi(z,\gamma)$ are defined in \eqref{eq:backstepping_z} and \eqref{eq:psi_definition}, respectively. Let $\Delta$ be given by either
\begin{align}
\Delta &= \delta, \label{eq:Delta_globa} \\
\Delta &= 2 \tan \frac{\delta}{2}. \label{eq:Delta_Barfli}
\end{align}
Then, the origin $\rho = \delta = \gamma = 0$ is globally asymptotically stable on $\mathcal{S}$ in the case of \eqref{eq:Delta_globa}, and on $\mathcal{S}_1$ in the case of \eqref{eq:Delta_Barfli}, in the sense of Definition~\ref{def-our-GAS}.
Furthermore, all the composite Lyapunov functions $V(\rho, \delta, \gamma) = \mathcal{V}(\rho^2, V_{\delta \gamma})$ and $V(\rho, \delta, \gamma) = \mathcal{V}(V_{\delta \gamma},\rho^2)$, for all functions $\mathcal{V}$ satisfying the conditions in Proposition \ref{prop:composite_Lyap_function}, and with $V_{\delta \gamma}$ for both cases of $\Delta$ defined as
\begin{equation}
\label{eq:V_CLF_backstepping}
V_{\delta \gamma}(\delta, \gamma) = 
\Delta^2 + q^2z^2\,, \quad q = \sqrt{\frac{k_1}{k_3}}\,,
\end{equation}
are (globally) strict CLFs for \eqref{eq:unicycle_polar_closed_loop-Gv-1} with respect to the input pair $(v/\rho,\omega)$ in the sense of Definition \ref{def-CLF}. 
\end{theorem}

\begin{lemma}\label{lemma:sinterm_upperbound}
The inequality $1- k \cos\gamma(1+\cos\gamma) \leq 2 (1 + k)  \tan^2\frac{\gamma}{2}$ holds for all $k\geq 1$ and $x\in\mathbb{R}$.
\end{lemma}

\begin{theorem}
\label{thm:unicycle_CLF_BoLSA}
Consider the system \eqref{eq:unicycle_polar_closed_loop-Gv-1} in closed-loop with \eqref{eq-basic-v-control}, \eqref{eq-omega-general} and 
\begin{equation}
\label{eq-control-bounded-in-gamma}
\tilde\omega =  k_2\sin\gamma + k_3 \displaystyle
% \cos\gamma \frac{(1+\cos\gamma)^2}{4}
\frac{\cos\gamma}{\left(1+\displaystyle\tan^2\frac{\gamma}{2}\right)^2}\delta \,, 
\end{equation}
with  $k_1, k_2, k_3 > 0$ such that 
$k_1 k_3\geq k_2^2$.
The origin $\rho = \delta = \gamma = 0$ is GAS on $\mathcal{S}_2$ in accordance with Definition \ref{def-our-GAS}.
Furthermore, all the composite Lyapunov functions $V(\rho, \delta, \gamma) = \mathcal{V}(\rho^2, V_{\delta \gamma})$ and $V(\rho, \delta, \gamma) = \mathcal{V}(V_{\delta \gamma},\rho^2)$, for all functions $\mathcal{V}$ satisfying the conditions in Proposition \ref{prop:composite_Lyap_function}, and with $V_{\delta \gamma}$ defined as
\begin{equation}
\hspace*{-0.2cm}
\label{eq:CLF_Bolsa_delta_gamma}
V_{\delta \gamma}(\delta, \gamma) =  k_3  \left(1+
%\frac{q}{k_2}\right) U +\frac{k_3}{2qk_2}U^2
\frac{2q^2+U}{2qk_2}\right)U
+   \left (\delta+2q\tan \frac{\gamma}{2} \right)^2\,,
\end{equation}
where
\begin{align}
U=\delta^2+  4q^2 \tan^2\frac{\gamma}{2} \,, \quad q = \sqrt{\frac{k_1}{k_3}}
\end{align} 
are (globally) strict CLFs for \eqref{eq:unicycle_polar_closed_loop-Gv-1} with respect to the input pair $(v/\rho,\omega)$ in the sense of Definition \ref{def-CLF}. 
\end{theorem}

\begin{proof}
We start by developing the feedback law \eqref{eq-control-bounded-in-gamma} and the CLF \eqref{eq:CLF_Bolsa_delta_gamma} for the $(\delta, \gamma)$-subsystem   \eqref{eq:unicycle_polar_closed_loop-Gv-3-pass}. 
% For this, consider the trigonometric indentities:
% \begin{eqnarray}
% \tan\frac{\gamma}{2} &=& \frac{\sin\gamma}{1+\cos\gamma} \label{eq:trig_identity_1}
% \\
% \left(\tan\frac{\gamma}{2}\right)^\prime &=& \frac{1}{1+\cos\gamma}  = \frac{1}{2}\left(1+\tan^2\frac{\gamma}{2}\right) \label{eq:trig_identity_2}
% \\
% \left(\tan^2\frac{\gamma}{2}\right)' &=& \frac{2 }{\sin \gamma} \tan^2 \frac{\gamma}{2} =
% 2 \frac{\sin\gamma}{(1+\cos\gamma)^2} \label{eq:trig_identity_3}
% \end{eqnarray}
We are going to use the following Lyapunov terms:
\begin{eqnarray}
U &=& \delta^2 +   q^2 V_0, \quad V_0 = 4 \tan^2\frac{\gamma}{2}\,,  \label{eq:U_bolsa} \\
% \\
% \Omega_1 &=&  V_0 + \frac{(\sqrt{k_3}\delta)^2}{2}, \quad k_3 > 0 \\
\Pi_0 &=& z^2\,, \quad z = \delta + 2 q \tan\frac{\gamma}{2}.
\end{eqnarray}
% \begin{equation}
% \fbox{$\displaystyle z = \delta + 2 q \tan\frac{\gamma}{2}, \quad q = \sqrt{\frac{k_1}{k_3}}$}, 
% \end{equation}
% Note that $V_0(\gamma)$ is radially unbounded on the interval $\gamma \in (-\pi,\pi)$, namely, it disallows ``winding'' the LOS angle to arbitrary real values. 
% It also disallows an initial condition facing away from the target, $\gamma=\pm\pi$, but this initial condition is neither disallowed by the forward velocity feedback nor by the angular velocity feedback we design next. 
The time derivative of the Lyapunov expression \eqref{eq:U_bolsa} along the solutions of \eqref{eq:unicycle_polar_closed_loop-Gv-3-pass} is
\begin{equation}
\dot U  =  \frac{8 q^2 \sin\gamma}{(1+\cos\gamma)^2} \left[ k_3 \delta \frac{(1+\cos\gamma)^2}{4} \cos\gamma 
- \tilde{\omega}\right]\,.
\end{equation}
We pick the control as in \eqref{eq-control-bounded-in-gamma}, where we use the fact that $(1+\cos\gamma)^2/4 = 1/(1+\tan^2 \gamma/2)^2$. With this feedback, we get
\begin{equation}
% \dot U  = -2 \frac{k_2  q^2 4 \sin^2\gamma}{(1+\cos\gamma)^2} = - 2 k_2  q^2 4 \tan^2\frac{\gamma}{2}  = - 2 k_2 q^2 V_0\,,
\dot U  = -2 \frac{k_2  q^2 4 \sin^2\gamma}{(1+\cos\gamma)^2} = - 2 k_2 q^2 V_0\,,
\end{equation}
and the dynamics for the LoS angle as
\begin{equation}
\dot\gamma = - k_2 \sin\gamma -  k_3 \delta\frac{(1+\cos\gamma)^2}{4} \cos\gamma  \,.
\end{equation}
Next, we consider the dynamics of the error $z$, which with
 $\sin(\gamma)\cos(\gamma) = \cos\gamma(1+\cos\gamma) \tan\frac{\gamma}{2}$ and $\frac{\sin\gamma}{1+\cos\gamma} =\tan\frac{\gamma}{2} $, are given by
\begin{equation}
\hspace*{-0.3cm}
\dot z = -k_2z + k_2\delta + k_3 q \frac{\cos\gamma(1+\cos\gamma)}{2}\left( 2 q \tan\frac{\gamma}{2}-\delta\right)\,.    
\end{equation}
Then,
\begin{equation}
\hspace*{-0.2cm}
\begin{aligned}[b]
\dot \Pi_0 = - 2k_2z^2 + 2k_2 z \delta &+ 
4 k_1 q \cos\gamma(1+\cos\gamma) \tan^2\frac{\gamma}{2} \\
&  
-k_3 q \frac{\cos\gamma(1+\cos\gamma)} {2}\delta^2\,.
\end{aligned}
\end{equation}
With $z\delta \leq  \frac{z^2}{4} + \delta^2$ and $\cos\gamma(1+\cos\gamma) \tan^2\frac{\gamma}{2} \leq  \frac{V_0}{2}
$, we get
\begin{equation}
\begin{aligned}[b]
    \dot \Pi_0 \le - &\frac{3}{2}k_2 \Pi_0 + 2 k_ 1 q V_0 \\
&+ 2 k_2 \delta^2\left(1- q \frac{k_3 }{2 k_2} 
\cos\gamma(1+\cos\gamma)\right)\,.
\end{aligned}
\end{equation}
It then follows from Lemma~\ref{lemma:sinterm_upperbound} with $k_1 k_3 \ge  k_2^2$ that 
\begin{equation}
\dot\Pi_0 
% &\leq  - \frac{3}{2}\Pi_0 +2 k_1 q V_0 +4 k_3 q\delta^2  \tan^2\frac{\gamma}{2} \nonumber\\
% &
\leq - \frac{3}{2} k_2 \Pi_0 + 2 k_1 q  V_0 +   2 k_3 q \delta^2 V_0.
\end{equation}
Taking into account that $\frac{k_3 }{k_2} q \dot U = - 2k_2 q^2 V_0$, and $\frac{k_3}{2q k_2 } \dot{U}^2 = -2k_3 q \delta^2 V_0 - 2k_1 q V_0^2$ and denoting
\begin{equation}
\Pi_1 = \Pi_0 + \frac{k_3 }{k_2} q  U  + \frac{k_3 }{2 q k_2} U^2\,, \label{eq:Pi_1_bolsa}
\end{equation}
we get
\begin{equation}
\dot\Pi_1 \leq - \frac{3}{2}k_2\Pi_0 - 2 k_1 q V_0^2\,.
\end{equation}
In conclusion, from \eqref{eq:U_bolsa} and \eqref{eq:Pi_1_bolsa}, we construct the Lyapunov function for the $(\delta, \gamma)$-subsystem as $V_{\delta \gamma} = k_3 U + \Pi_1$ which is equivalent to \eqref{eq:CLF_Bolsa_delta_gamma} 
% \begin{equation}
% V_{\delta \gamma} (\delta, \gamma) = k_3 U  + \Pi_1 =  k_3  \left(1+
% %\frac{q}{k_2}\right) U +\frac{k_3}{2qk_2}U^2
% \frac{2q^2+U}{2qk_2}\right)U
% +   \left (\delta+2q\tan \frac{\gamma}{2} \right)^2
% %\nonumber\\ &=& \frac{\rho^2}{2} + 
% \,, \label{eq:tan_Lyapunov_function}
% \end{equation}
and has the time derivative \begin{equation}
\dot V_{\delta \gamma} \le  - 2k_1 k_2 V_0 - \frac{3}{2}k_2 \left(\delta + q \tan \frac{\gamma}{2}\right)^2 - 2k_1 q V_0^2
\label{eq:V_dot_Bolsa}
\end{equation} 
which is negative definite on $\mathcal{T}_2$. Analogous to the proof of Theorem~\ref{thm:backstepping_theorem}, we arrive at both statements of Theorem~\ref{thm:unicycle_CLF_BoLSA}.
% conclude that the composite Lyapunov functions $V(\rho, \delta, \gamma) = \mathcal{V}(\rho^2, V_{\delta \gamma})$ and $V(\rho, \delta, \gamma) = \mathcal{V}(V_{\delta \gamma},\rho^2)$ with  $\dot{V}(\rho,\delta,\gamma) \le - \alpha(V)$ are strict CLFs for \eqref{eq:unicycle_polar_closed_loop-Gv-1}. 
% We also have that 
% $\abs{(\rho,\delta,\gamma)}_{\mathcal{S}_2} \le \bar{\alpha}_1^{-1}(\beta( \bar{\alpha}_2(\abs{(\rho_0, \delta_0,\gamma_0)}_{\mathcal{S}_2}),t))$
% where $\bar{\alpha}_1^{-1}(\beta( \bar{\alpha}_2(r),t))$ is class $\mathcal{KL}$ in $(r,t)$.
% Thus, 
% %we  conclude that the origin 
% $\rho =  \delta = \gamma = 0$ is GAS on $\mathcal{S}_2$.
\end{proof}

\end{document}
